\documentclass[12pt,a4paper]{article}

\usepackage[spanish,es-nodecimaldot]{babel}
\usepackage[utf8]{inputenc}
\usepackage[T1]{fontenc}
\usepackage{lmodern}
\usepackage{geometry}
\usepackage{microtype}
\usepackage{setspace}
\usepackage{parskip}
\usepackage{enumitem}
\usepackage{xcolor}
\usepackage{listings}
\usepackage{hyperref}
\usepackage{amsmath}
\usepackage{amssymb}
\usepackage{array}
\usepackage{booktabs}
\usepackage{longtable}

\geometry{margin=2.5cm}
\onehalfspacing
\setlist[itemize]{leftmargin=*,itemsep=0.2em}
\setlist[enumerate]{leftmargin=*,itemsep=0.2em}

\hypersetup{
  colorlinks=true,
  linkcolor=black,
  urlcolor=blue,
  pdftitle={Resolucion del Practico 7 - Testing de Software 2024},
  pdfauthor={}
}

\lstdefinestyle{codestyle}{
  basicstyle=\ttfamily\small,
  frame=single,
  breaklines=true,
  showstringspaces=false,
  keywordstyle=\color{blue!70!black}\bfseries,
  commentstyle=\color{gray!70!black}\itshape,
  stringstyle=\color{orange!70!black}
}

\title{\textbf{Resolucion del Practico 7}\\Testing de Software 2024}
\author{}
\date{}

\begin{document}

\maketitle
\tableofcontents
\newpage

\section*{Introduccion}
\addcontentsline{toc}{section}{Introduccion}

Este documento contiene la resolucion del \textit{Practico 7} de Testing de Software, centrado
en \textbf{Property-Based Testing (PBT)} usando la libreria \texttt{jqwik} sobre Java. En lugar
de escribir tests con valores concretos, se enuncian \emph{propiedades} que deben valer para
\emph{todas} las entradas dentro de un dominio razonable, y el framework las verifica generando
cientos de instancias aleatorias.

A lo largo del practico se trabajan tres ejercicios:

\begin{itemize}
  \item \textbf{Ejercicio 1.} \texttt{NodeCachingLinkedList}: una lista doblemente enlazada y
  circular con un \emph{pool} de nodos cacheados. Hay que completar \texttt{repOK()} con todos
  los invariantes de representacion y escribir tres propiedades \texttt{jqwik} que validen el
  comportamiento.
  \item \textbf{Ejercicio 2.} \texttt{Point}: clase de puntos con coordenadas \texttt{float}.
  Se implementan \texttt{equals}/\texttt{hashCode} de manera consistente con el contrato de
  \texttt{Object}, se escribe un generador \texttt{@Provide} para puntos y dos propiedades:
  el contrato \texttt{equals/hashCode} y la distancia entre puntos sobre una recta paralela al
  eje X.
  \item \textbf{Ejercicio 3.} \texttt{Date}: representacion de fechas a partir de 1900. Se
  completa el constructor, \texttt{repOk()} y el metodo \texttt{addDays}, y se enuncia la
  propiedad \emph{toda fecha valida sumada a un numero no negativo de dias devuelve una fecha
  valida}.
\end{itemize}

\textbf{Convenciones.} A lo largo del documento se usa la notacion estandar de \texttt{jqwik}:
\texttt{@Property} marca metodos como propiedades, \texttt{@ForAll} parametros sobre los que
\texttt{jqwik} muestrea valores, y \texttt{@Provide} define generadores. Cada propiedad se
ejecuta con \texttt{tries = 200} o \texttt{250} iteraciones segun el caso, equilibrando
cobertura y tiempo de ejecucion.

\textbf{Como se corren los tests.} Desde \texttt{tp7/assignmnet-7-rodeghiero}:

\begin{lstlisting}[style=codestyle]
JAVA_HOME=$(/usr/libexec/java_home -v 17) mvn -Dmaven.repo.local=.m2 \
    -Djacoco.skip=true test
\end{lstlisting}

La flag \texttt{-Djacoco.skip=true} se usa porque el plugin de cobertura del template
(\texttt{0.8.2}) es incompatible con la JVM mas reciente; PBT no necesita ese instrumento.

\section{Ejercicio 1: \texttt{NodeCachingLinkedList}}

\subsection{Estructura del objeto}

La clase mantiene dos sub-estructuras que conviven:

\begin{itemize}
  \item \textbf{Lista principal}: doblemente enlazada y \emph{circular}, con un nodo centinela
  \texttt{header}. Cada nodo \texttt{n} cumple
  $\texttt{n.next.previous} = \texttt{n}$ y $\texttt{n.previous.next} = \texttt{n}$.
  \item \textbf{Cache}: una lista \emph{simplemente} enlazada, accesible por
  \texttt{firstCachedNode}, que mantiene nodos reciclados. Cuando se remueve un elemento, el
  nodo se ``limpia'' (\texttt{previous = null}, \texttt{value = null}) y se devuelve a la
  cache hasta un maximo de \texttt{maximumCacheSize}.
\end{itemize}

\subsection{(a) Implementacion de \texttt{repOK()}}

Los invariantes son los siguientes, agrupados por costo creciente de chequeo:

\begin{itemize}
  \item Validez basica: \texttt{header != null}, \texttt{header.next != null},
  \texttt{header.previous != null}.
  \item Consistencia de tamanos: \texttt{size >= 0}, \texttt{cacheSize >= 0},
  \texttt{maximumCacheSize >= 0} y \texttt{cacheSize <= maximumCacheSize}.
  \item Constante: \texttt{DEFAULT\_MAXIMUM\_CACHE\_SIZE == 20}, como pide el enunciado.
  \item Lista principal: recorrer desde \texttt{header} siguiendo \texttt{next} hasta volver al
  centinela. En cada paso verificar los enlaces dobles y, con un \texttt{HashSet}, que no haya
  ciclos prematuros. Al final, \texttt{size} debe coincidir con el numero de nodos visitados
  menos uno (por el centinela).
  \item Cache: recorrer desde \texttt{firstCachedNode} siguiendo \texttt{next} hasta
  \texttt{null}. Cada nodo cacheado debe tener \texttt{previous == null} y
  \texttt{value == null}, no formar ciclos y \emph{no} pertenecer tambien a la lista principal
  (se compara contra el \texttt{HashSet} anterior).
  \item Al final \texttt{cacheSize} debe coincidir con la cantidad de nodos recorridos en la
  cache.
\end{itemize}

\textbf{Resumen del codigo.}

\begin{lstlisting}[style=codestyle,language=Java]
public boolean repOK() {
    if (header == null) return false;
    if (header.getNext() == null || header.getPrevious() == null) return false;
    if (size < 0) return false;
    if (maximumCacheSize < 0 || cacheSize < 0 ||
        cacheSize > maximumCacheSize) return false;
    if (DEFAULT_MAXIMUM_CACHE_SIZE != 20) return false;

    Set<LinkedListNode> listNodes = new HashSet<>();
    LinkedListNode current = header;
    int listCount = 0;
    while (true) {
        if (current == null) return false;
        if (!listNodes.add(current)) return false;
        if (current.getNext() == null || current.getPrevious() == null) return false;
        if (current.getNext().getPrevious() != current) return false;
        if (current.getPrevious().getNext() != current) return false;
        listCount++;
        current = current.getNext();
        if (current == header) break;
    }
    if (size != listCount - 1) return false;

    Set<LinkedListNode> cacheNodes = new HashSet<>();
    current = firstCachedNode;
    int cacheCount = 0;
    while (current != null) {
        if (!cacheNodes.add(current)) return false;
        if (listNodes.contains(current)) return false;
        if (current.getPrevious() != null) return false;
        if (current.getValue() != null) return false;
        cacheCount++;
        current = current.getNext();
    }
    return cacheSize == cacheCount;
}
\end{lstlisting}

\subsection{(b) Inicializacion de \texttt{maximumCacheSize}}

El constructor original dejaba \texttt{maximumCacheSize} en $0$, lo cual desactiva la cache
(porque \texttt{isCacheFull()} devuelve siempre \texttt{true} y ningun nodo se cachea). Eso es
inconsistente con el nombre de la clase. Se corrige asignando
\texttt{maximumCacheSize = DEFAULT\_MAXIMUM\_CACHE\_SIZE} en el constructor:

\begin{lstlisting}[style=codestyle,language=Java]
public NodeCachingLinkedList() {
    header = new LinkedListNode();
    header.setValue(null);
    header.setNext(header);
    header.setPrevious(header);
    maximumCacheSize = DEFAULT_MAXIMUM_CACHE_SIZE;   // <-- agregado
}
\end{lstlisting}

\subsection{(c) Propiedades \texttt{jqwik}}

\paragraph{Propiedad 1: al remover, la cache crece en uno.}

\begin{lstlisting}[style=codestyle,language=Java]
@Property(tries = 200)
void luegoDeRemoverUnElementoSeIncrementaEnUnoElTamanoDeCache(
    @ForAll("escenariosRemocion") EscenarioRemocion escenario) {
    NodeCachingLinkedList ncl = escenario.lista;
    assertTrue(ncl.repOK());
    int cacheAntes = ncl.getCacheSize();

    Integer removido = ncl.removeIndex(escenario.index);

    assertNotNull(removido);
    assertEquals(cacheAntes + 1, ncl.getCacheSize());
}
\end{lstlisting}

El generador construye una lista de entre 1 y 15 elementos. Como ese rango es menor a
\texttt{maximumCacheSize = 20}, la cache nunca llega a saturarse y la propiedad ``cada
remocion suma 1 al \texttt{cacheSize}'' vale siempre.

\paragraph{Propiedad 2: con cache no vacia, \texttt{size + cacheSize} se conserva en
\texttt{addFirst}.}

\begin{lstlisting}[style=codestyle,language=Java]
@Property(tries = 200)
void siLaCacheNoEstaVaciaAgregarMantieneConstanteLaSumaDeNodos(
    @ForAll("escenariosAgregarConCacheNoVacia") EscenarioAgregarConCache escenario) {
    NodeCachingLinkedList ncl = escenario.lista;
    assertTrue(ncl.repOK());
    assertTrue(ncl.getCacheSize() > 0);
    int sumaAntes = ncl.getSize() + ncl.getCacheSize();

    ncl.addFirst(escenario.nuevoValor);

    int sumaDespues = ncl.getSize() + ncl.getCacheSize();
    assertEquals(sumaAntes, sumaDespues);
}
\end{lstlisting}

Esta propiedad captura la idea central de la cache: si hay nodos disponibles, \texttt{addFirst}
\emph{reutiliza} un nodo en lugar de crear uno. Por lo tanto, \texttt{size} aumenta en 1 y
\texttt{cacheSize} disminuye en 1, manteniendo constante la suma. Para garantizar la
precondicion (cache no vacia), el generador primero construye una lista y despues remueve un
elemento.

\paragraph{Propiedad 3: \texttt{repOK()} se preserva por \texttt{removeIndex}.}

\begin{lstlisting}[style=codestyle,language=Java]
@Property(tries = 200)
void eliminarUnElementoMantieneElInvarianteDeRepresentacion(
    @ForAll("escenariosRemocion") EscenarioRemocion escenario) {
    NodeCachingLinkedList ncl = escenario.lista;
    assertTrue(ncl.repOK());
    ncl.removeIndex(escenario.index);
    assertTrue(ncl.repOK());
}
\end{lstlisting}

Esta es la propiedad clasica de un invariante de representacion: cualquier operacion publica
sobre un objeto valido debe dejarlo en estado valido. Si alguna manipulacion interna (manejo
de enlaces, paso del nodo a la cache, decremento de \texttt{size}) rompiera la coherencia, el
segundo \texttt{repOK()} la detectaria.

\subsection{Generadores}

Los \texttt{@Provide} construyen escenarios completos para que cada propiedad reciba
exactamente lo que necesita:

\begin{lstlisting}[style=codestyle,language=Java]
@Provide
Arbitrary<EscenarioRemocion> escenariosRemocion() {
    return Arbitraries.integers().between(1, 15).flatMap(size ->
        Arbitraries.integers().between(-1000, 1000).list().ofSize(size).flatMap(valores ->
            Arbitraries.integers().between(0, size - 1).map(index ->
                new EscenarioRemocion(crearLista(valores), index))));
}

@Provide
Arbitrary<EscenarioAgregarConCache> escenariosAgregarConCacheNoVacia() {
    return Arbitraries.integers().between(1, 15).flatMap(size ->
        Arbitraries.integers().between(-1000, 1000).list().ofSize(size).flatMap(valores ->
            Arbitraries.integers().between(0, size - 1).flatMap(indexARemover ->
                Arbitraries.integers().between(-1000, 1000).map(nuevoValor -> {
                    NodeCachingLinkedList ncl = crearLista(valores);
                    ncl.removeIndex(indexARemover);
                    return new EscenarioAgregarConCache(ncl, nuevoValor);
                }))));
}
\end{lstlisting}

El uso de \texttt{flatMap} es necesario porque el rango del indice depende del tamano de la
lista: \texttt{Arbitraries.integers().between(0, size - 1)} solo se puede armar despues de
haber generado \texttt{size}.

\section{Ejercicio 2: \texttt{Point}}

\subsection{Contrato \texttt{equals/hashCode}}

El contrato definido por \texttt{Object} es:

\begin{itemize}
  \item Si \texttt{a.equals(b)}, entonces \texttt{a.hashCode() == b.hashCode()}. \emph{Esta es
  la regla obligatoria}: si se viola, las estructuras basadas en hashing (\texttt{HashMap},
  \texttt{HashSet}) dejan de funcionar correctamente.
  \item La inversa no es obligatoria: dos objetos pueden tener el mismo \texttt{hashCode} sin
  ser iguales (\emph{colision}); es inevitable porque el hash es un \texttt{int} con rango
  finito.
\end{itemize}

\subsection{Implementacion}

\begin{lstlisting}[style=codestyle,language=Java]
@Override
public boolean equals(Object obj) {
    if (this == obj) return true;
    if (obj == null || getClass() != obj.getClass()) return false;
    Point other = (Point) obj;
    return Float.compare(this.x, other.x) == 0
        && Float.compare(this.y, other.y) == 0;
}

@Override
public int hashCode() {
    int result = Float.floatToIntBits(x);
    result = 31 * result + Float.floatToIntBits(y);
    return result;
}
\end{lstlisting}

\textbf{Por que \texttt{Float.compare} en \texttt{equals}.} El operador \texttt{==} tiene
problemas con los valores especiales del estandar IEEE 754: \texttt{NaN == NaN} es
\texttt{false}, lo que romperia la reflexividad de \texttt{equals}. \texttt{Float.compare}
trata correctamente \texttt{NaN} y la distincion entre \texttt{-0.0f} y \texttt{+0.0f}.

\textbf{Por que \texttt{Float.floatToIntBits} en \texttt{hashCode}.} Convierte el \texttt{float}
a su representacion binaria como \texttt{int}. Asi, dos \texttt{float} iguales segun
\texttt{Float.compare} producen el mismo \texttt{int} (y por lo tanto el mismo hash). La
combinacion \texttt{31 * result + ...} es la formula clasica para hashear campos multiples
con buena distribucion.

\subsection{Propiedad 1: contrato \texttt{equals/hashCode}}

\begin{lstlisting}[style=codestyle,language=Java]
@Property(tries = 250)
void puntosIgualesDebenTenerMismoHashCode(@ForAll("puntos") Point punto) {
    Point mismoPunto = new Point(punto.getX(), punto.getY());
    assertTrue(punto.equals(mismoPunto));
    assertEquals(punto.hashCode(), mismoPunto.hashCode());
}

@Provide
Arbitrary<Point> puntos() {
    return Combinators.combine(coordenadas(), coordenadas()).as(Point::new);
}

@Provide
Arbitrary<Float> coordenadas() {
    return Arbitraries.floats().between(-10_000f, 10_000f);
}
\end{lstlisting}

El generador construye puntos con coordenadas en $[-10^4, 10^4]$, rango razonable para
ejercitar la combinacion de bits sin caer en valores patologicos como \texttt{NaN} o
$\pm\infty$.

\subsection{Propiedad 2: distancia sobre recta horizontal}

Si $p_1 = (x_1, y)$ y $p_2 = (x_2, y)$, entonces $\textit{distance}(p_1, p_2) = |x_2 - x_1|$,
porque la ordenada se cancela en la formula euclidea. Esta es una propiedad geometrica
\emph{independiente} de los valores concretos: vale para todo $x_1, x_2, y$ finitos.

\begin{lstlisting}[style=codestyle,language=Java]
@Property(tries = 250)
void distanciaEnRectaParalelaAlEjeXEsDiferenciaDeAbscisas(
    @ForAll("coordenadas") float x1,
    @ForAll("coordenadas") float x2,
    @ForAll("coordenadas") float y) {
    Point p1 = new Point(x1, y);
    Point p2 = new Point(x2, y);
    double distanciaEsperada = Math.abs((double) (x2 - x1));
    double distanciaReal = p1.distanceTo(p2);
    assertEquals(distanciaEsperada, distanciaReal, 1.0e-6);
}
\end{lstlisting}

\textbf{Sobre la tolerancia.} \texttt{distanceTo} usa \texttt{Math.sqrt(Math.pow(...))}, que
introduce errores de redondeo. Por eso se usa \texttt{assertEquals} con tolerancia $10^{-6}$.
Sin tolerancia, la propiedad fallaria por ruido numerico aun siendo logicamente correcta.

\section{Ejercicio 3: \texttt{Date}}

\subsection{Constructor y \texttt{repOk()}}

El constructor valida la fecha antes de inicializar los campos, y \texttt{repOk()} reusa la
misma logica:

\begin{lstlisting}[style=codestyle,language=Java]
public Date(int d, int m, int y) throws IllegalArgumentException {
    if (!isValidDate(d, m, y)) throw new IllegalArgumentException("invalid date");
    this.day = d; this.month = m; this.year = y;
    assert repOk();
}

public boolean repOk() { return isValidDate(day, month, year); }

private static boolean isValidDate(int d, int m, int y) {
    if (y < 1900) return false;
    if (m < 1 || m > 12) return false;
    if (d < 1) return false;
    return d <= daysInMonth(m, y);
}

private static int daysInMonth(int m, int y) {
    switch (m) {
        case 1: case 3: case 5: case 7: case 8: case 10: case 12: return 31;
        case 4: case 6: case 9: case 11: return 30;
        case 2: return leap(y) ? 29 : 28;
        default: return -1;
    }
}
\end{lstlisting}

\textbf{Por que reusar.} Tener \emph{una sola} implementacion de \texttt{daysInMonth} y de la
validacion evita que constructor y \texttt{addDays} se desincronicen. Cualquier cambio
(eg.\ admitir anos $< 1900$) se hace en un solo lugar.

\subsection{\texttt{addDays(when, days)}}

\textbf{Estrategia:} avanzar por bloques de mes en vez de dia por dia. Para \texttt{days} del
orden de miles, una iteracion por dia seria innecesariamente lenta.

\begin{lstlisting}[style=codestyle,language=Java]
public Date addDays(Date when, int days) {
    if (when == null) throw new IllegalArgumentException("date cannot be null");
    if (days < 0)     throw new IllegalArgumentException("days must be non-negative");

    int d = when.day, m = when.month, y = when.year;
    int remainingDays = days;

    while (remainingDays > 0) {
        int daysInCurrentMonth = daysInMonth(m, y);
        int daysLeftInMonth    = daysInCurrentMonth - d;

        if (remainingDays <= daysLeftInMonth) {
            d = d + remainingDays;
            remainingDays = 0;
        } else {
            remainingDays = remainingDays - (daysLeftInMonth + 1);
            d = 1;
            m = m + 1;
            if (m > 12) { m = 1; y = y + 1; }
        }
    }
    return new Date(d, m, y);
}
\end{lstlisting}

\textbf{Invariantes del algoritmo:}
\begin{itemize}
  \item \texttt{(d, m, y)} siempre representa una fecha \emph{valida}.
  \item Despues de cada iteracion, la suma ``dias ya consumidos $+$ \texttt{remainingDays}''
  es igual a \texttt{days}.
\end{itemize}

\textbf{El \texttt{+ 1} en el caso ``no cabe''.} Cuando \texttt{remainingDays > daysLeftInMonth}
se consumen \texttt{daysLeftInMonth} dias para llegar al ultimo dia del mes y \emph{uno mas}
para pasar al dia 1 del mes siguiente; de ahi \texttt{(daysLeftInMonth + 1)}.

\textbf{Por que reconstruir con \texttt{new Date(d, m, y)} al final.} El constructor valida.
Si por algun bug la fecha intermedia quedara mal, la excepcion lo expondria de inmediato.

\subsection{Propiedad: \texttt{addDays} siempre devuelve una fecha valida}

\begin{lstlisting}[style=codestyle,language=Java]
@Property(tries = 250)
void addDaysSiempreDevuelveFechaValida(
    @ForAll("fechasValidas") Date fecha,
    @ForAll("diasNoNegativos") int dias) {
    Date receptor = new Date(1, 1, 1900);
    Date resultado = receptor.addDays(fecha, dias);
    assertNotNull(resultado);
    assertTrue(resultado.repOk());
}
\end{lstlisting}

\paragraph{Generadores.}

\begin{itemize}
  \item \texttt{fechasValidas} (\texttt{@Provide}) usa \texttt{flatMap} para que el rango del
  dia dependa del mes y del ano:
  \begin{lstlisting}[style=codestyle,language=Java]
  Arbitraries.integers().between(1900, 2400).flatMap(year ->
    Arbitraries.integers().between(1, 12).flatMap(month ->
      Arbitraries.integers().between(1, daysInMonth(month, year))
        .map(day -> new Date(day, month, year))));
  \end{lstlisting}
  Sin este encadenamiento se generarian cosas como \texttt{31/04/2024}, que el constructor
  rechazaria.
  \item \texttt{diasNoNegativos}: enteros en $[0, 5000]$. El rango es lo suficientemente amplio
  como para forzar saltos de varios anos (5000 dias $\approx$ 13.7 anos) sin volver la suite
  lenta.
\end{itemize}

\paragraph{Sobre el receptor.} \texttt{addDays} es un metodo de instancia. Se llama sobre un
\texttt{new Date(1, 1, 1900)} arbitrario; el receptor no se usa en el calculo (todo se basa en
\texttt{when}), asi que cualquier instancia valida sirve para invocarlo.

\section*{Conclusiones}
\addcontentsline{toc}{section}{Conclusiones}

\begin{itemize}
  \item El testing basado en propiedades cambia el foco: en vez de pensar en
  \emph{ejemplos}, se enuncia el comportamiento \emph{en general}. El framework se encarga de
  buscar entradas que rompan la propiedad.
  \item Las propiedades mas utiles son las que capturan \emph{invariantes} (e.g.\ ``\texttt{repOK}
  se preserva''), \emph{contratos} (\texttt{equals/hashCode}) y \emph{relaciones algebraicas o
  geometricas} (distancia sobre recta horizontal).
  \item Los \textbf{generadores} son tan importantes como las propiedades: tienen que producir
  entradas validas y representativas. \texttt{flatMap} permite construir generadores
  \emph{dependientes}, donde un parametro depende de otro (e.g.\ dia $\le$ dias del mes).
  \item Un buen \texttt{repOK} no es solo documentacion: \emph{ejecuta} el invariante y lo
  hace verificable por las propiedades.
\end{itemize}

\end{document}
